\section{Wnioski końcowe}
\label{sec:wnioski}

\subsection{Podsumowanie}

W pracy zrealizowano empiryczne porównanie metod uczenia głębokiego i lasów
losowych w zadaniu prognozowania naruszeń danych w chmurze, sformułowanym jako
binarna klasyfikacja przepływów sieciowych. Zbudowano kompletny, odtwarzalny
potok obejmujący przygotowanie danych, trening pięciu modeli (regresja
logistyczna, las losowy oraz sieci MLP, CNN 1D i LSTM) oraz ich ocenę pod kątem
skuteczności i kosztu obliczeniowego.

Najważniejsze wnioski są następujące:
\begin{enumerate}
  \item \textbf{Obie rodziny metod skutecznie wykrywają naruszenia.} Zarówno las
        losowy, jak i sieci głębokie osiągają wysokie wyniki, wyraźnie
        przewyższając prosty model liniowy. To pokazuje, że złośliwa aktywność
        kryje się w złożonych zależnościach między cechami, które obie metody
        dobrze wychwytują.
  \item \textbf{Skuteczność jest niemal taka sama, ale koszt --- już nie.}
        Najlepszy model głęboki (\resBestDeepModel, $F_1 = \resBestDeepFOne$) i
        las losowy ($F_1 = \resRFFOne$) osiągają praktycznie identyczną jakość,
        przy czym las losowy dodatkowo wykrywa najwięcej ataków (najwyższa
        czułość), uczy się kilkadziesiąt razy szybciej i potrafi wskazać, które
        cechy zaważyły na decyzji. To czyni go wyjątkowo atrakcyjnym wyborem.
  \item \textbf{Bardziej skomplikowane nie znaczy lepsze.} Rozbudowane sieci
        sekwencyjne (zwłaszcza LSTM) nie przewyższyły ani lasu losowego, ani
        prostszej sieci, generując jedynie wyższy koszt. Potwierdza to znaną z
        literatury obserwację, że dla danych w postaci tabeli metody oparte na
        drzewach pozostają bardzo trudne do pobicia.
  \item \textbf{Wybór metody zależy od kontekstu.} Jeśli liczą się szybkie
        dotrenowywanie, niski koszt, dobra wykrywalność i zrozumiałość decyzji
        --- las losowy jest naturalnym wyborem. Sieci głębokie warto rozważyć
        tam, gdzie zależy nam na maksymalnym ograniczeniu fałszywych alarmów
        (wysoka precyzja CNN) i dysponujemy większymi zasobami obliczeniowymi.
\end{enumerate}

\subsection{Ograniczenia}

Wyniki należy interpretować w świetle ograniczeń pracy:
\begin{itemize}
  \item Eksperyment przeprowadzono na losowym podzbiorze danych
        \texttt{\resDataSource} (ze względu na rozmiar pełnego zbioru liczącego
        kilka gigabajtów). Większa próbka mogłaby nieco zmienić konkretne
        liczby, choć raczej nie ogólne wnioski; program jest gotowy do
        uruchomienia na pełnym zbiorze bez zmian w kodzie.
  \item Nie przeprowadzono wyczerpującego strojenia hiperparametrów --- użyto
        rozsądnych wartości domyślnych, jednakowych dla wszystkich modeli.
  \item Obliczenia prowadzono na CPU; na akceleratorach GPU względne czasy
        treningu modeli głębokich uległyby skróceniu, choć relacja kosztów
        pozostałaby jakościowo podobna.
  \item Zadanie ograniczono do klasyfikacji binarnej; analiza wieloklasowa
        (rozróżnianie typów ataków) stanowi naturalne rozszerzenie.
\end{itemize}

\subsection{Kierunki dalszych prac}

Możliwe rozszerzenia obejmują: (1) walidację na pełnym zbiorze CSE-CIC-IDS2018 i
innych zbiorach (UNSW-NB15, NSL-KDD); (2) systematyczne strojenie
hiperparametrów oraz włączenie metod gradient boosting (XGBoost, LightGBM);
(3) klasyfikację wieloklasową i wykrywanie nieznanych ataków metodami
półnadzorowanymi; (4) analizę odporności modeli na ataki adwersarialne oraz
dryf danych; (5) integrację najlepszego modelu w architekturze strumieniowej
(np.\ przetwarzanie przepływów w czasie rzeczywistym) wraz z mechanizmami
wyjaśnialności (SHAP) dla wsparcia analityków SOC.

\subsection{Podsumowanie końcowe}

Przeprowadzone badanie pokazuje, że w prognozowaniu naruszeń danych w chmurze
nie ma jednej, bezdyskusyjnie najlepszej metody --- wybór jest kompromisem
między skutecznością a kosztem, szybkością i zrozumiałością decyzji. Dla danych
opisujących ruch sieciowy las losowy okazuje się wyjątkowo silnym i tanim
rozwiązaniem, któremu sieci głębokie co najwyżej dorównują, lecz go nie
prześcigają --- i to dopiero za cenę znacznie większego nakładu obliczeniowego.
Dla studenta czy inżyniera bezpieczeństwa płynie z tego praktyczny wniosek:
zanim sięgnie się po skomplikowane sieci neuronowe, warto najpierw wypróbować
dobrze dobraną metodę klasyczną.
